% Figure: a rectangular gate hinged at the bottom holding back water. The
% hydrostatic pressure grows linearly with depth, a triangular load whose
% resultant acts at the centre of pressure, one third of the depth up from
% the base. The moment about the hinge sets the latch force at the top.
\begin{figure}[htbp]
\centering
\begin{tikzpicture}[>=latex, x=1.0cm, y=1.0cm]
  % water region (left), shaded
  \fill[engblue!10] (-3.4,0) rectangle (0,3.2);
  \draw[engblue, thin] (-3.4,3.2) -- (0,3.2);
  \foreach \xx in {-3.2,-2.8,-2.4,-2.0,-1.6,-1.2,-0.8,-0.4} {
    \draw[engblue, thin] (\xx,3.35) -- (\xx+0.2,3.2);
  }
  \node[engblue, font=\scriptsize, anchor=south west] at (-3.3,3.35) {water surface};
  % the gate (vertical wall from base to surface)
  \draw[engblue, very thick] (0,0) -- (0,3.2);
  \node[engblue, font=\scriptsize, anchor=west, rotate=90] at (0.12,1.2) {gate, height $h$};
  % hinge at the base
  \draw[engblack, fill=enggray!30] (0,0) circle (0.12);
  \node[font=\scriptsize, anchor=north] at (0,-0.15) {hinge $O$};
  % latch at the top
  \draw[->, very thick, engorange] (0.9,3.2) -- (0.12,3.2);
  \node[engorange, font=\scriptsize, anchor=west] at (0.95,3.2) {$F$ (latch)};
  % triangular pressure diagram on the water side
  \draw[engred, thick] (0,3.2) -- (0,0) -- (-1.6,0) -- cycle;
  \foreach \yy/\len in {0.3/-0.15,0.7/-0.35,1.1/-0.55,1.5/-0.75,1.9/-0.95,2.3/-1.15,2.7/-1.35} {
    \draw[->, engred, thin] (\len,\yy) -- (0,\yy);
  }
  \node[engred, font=\scriptsize, anchor=north east] at (-1.6,0) {$p_{\max}=\rho g h$};
  % resultant at one-third height
  \draw[->, very thick, engred] (-2.9,1.067) -- (-0.1,1.067);
  \node[engred, font=\scriptsize, anchor=south east] at (-2.0,1.15) {$R=\tfrac12\rho g h^2$};
  \draw[<->, enggray, thin] (0.55,0) -- (0.55,1.067);
  \node[enggray, font=\scriptsize, anchor=west] at (0.6,0.53) {$h/3$};
\end{tikzpicture}
\caption[Hydrostatic pressure on a gate]{A vertical gate hinged at its base
$O$ and latched at the top holds back water of depth $h$. The hydrostatic
pressure grows linearly with depth from zero at the surface to
$p_{\max}=\rho g h$ at the base, a triangular distribution whose resultant per
unit width is $R=\tfrac12\rho g h^2$ acting at the centre of pressure, one
third of the depth up from the base where the triangle's centroid sits.
Summing moments about the hinge sets the latch force $F$ at the top: the
resultant's short lever arm about the base, $h/3$, works against the latch's
full arm $h$, so the latch carries only a fraction of the resultant. The
centre of pressure sitting below the mid-depth, not at it, is the fact a gate
or dam analysis lives by.}
\label{fig:vol03ch02:hydrostatic-gate}
\end{figure}
